\documentclass{article}
\usepackage{amsmath}
\usepackage{amssymb}
\usepackage{amsthm}
\usepackage{enumitem}
\usepackage[a4paper, margin=1in]{geometry}


\begin{document}
\setlength{\parindent}{0pt}

\section{Geometria in $\mathbb{R}^3$}
\subsection{Piani nello spazio}
    Dato un punto $P_0 \in \mathbb{R}^3$ e un vettore $\mathbf{v} \in \mathbb{R}^3$, è possibile individuare un piano passante per $P_0$ ortogonale a $\mathbf{v}$:
    \[
        (P - P_0) \cdot \mathbf{v} = 0 \iff a(x - x_0) + b(y - y_0) + c(z - z_0) = 0 \iff ax + by + cz + d = 0
    \] 
    L'equazione cartesiana di un piano è quindi $\pi: ax + by + cz + d = 0$.

    \vspace{10pt}
    Dati invece tre punti $P_0, P_1, P_2$ giacenti sul piano è possibile determinarne l'equazione parametrica:
    \[
        (P - P_0) = \lambda(P_0 - P_1) + \mu(P_1 - P_2) \iff 
        \begin{cases}
            x - x_0 = \lambda(x_0 - x_1) + \mu(x_1 - x_2) \\
            y - y_0 \, = \lambda(y_0 - y_1) + \mu(y_1 - y_2) \\
            z - z_0 \, = \lambda(z_0 - z_1) + \mu(z_1 - z_2)
        \end{cases}  
    \]



\subsection{Retta nello spazio}
    Dato un punto $P_0$ e un vettore $\mathbf{v}$ direzione è possibile individuare una retta passante per il punto:
    \[
        r: (P - P_0) = t\mathbf{v} \iff 
        \begin{cases}
            x - x_0 = a t \\
            y - y_0 = b t \\
            z - z_0 = c t 
        \end{cases}
    \]

    Per poter individuare l'equazione cartesiana della retta in uno spazio tridimensionale sono necessiare 2 equazioni,
    in quanto la codimensione dello spazio vettoriale $W \in \mathbb{R}^3$ rispetto ad $\mathbb{R}^3$ è 2:
    \[
        \begin{cases}
            x = x_0 + a t \\
            y = y_0 + b t \\
            z = z_0 + c t 
        \end{cases}
        \, 
        \iff
        \quad
        \begin{cases}
            y - y_0 + b(\frac{x - x_0}{a}) = 0 \\
            z - z_0 + c(\frac{x - x_0}{a}) = 0 \\
        \end{cases}
        \, 
        \iff
        \quad
        \begin{cases}
            y + \frac{b}{a}x - \frac{b}{a}x_0 - y_0 = 0\\
            z + \frac{c}{a}x - \frac{c}{a}x_0 - y_0 = 0\\
        \end{cases}
    \]
    Di conseguenza, per individuare una retta nello spazio in forma cartesiana è necessario esprimerla come intersezione di due piani.



\subsection{Fasci di piani}
    Data una retta scritta in forma cartesiana è possibile individuare un fascio di piani passante per la retta $r$ mediante 
    una combinazione lineare dei piani generatori:
    \[
        r: \begin{cases} ax + by + cz + d = 0 \\ a'x + b'y + c'z + d' = 0\end{cases} \implies \quad \lambda(ax + by + cz + d) + \mu(a'x + b'y + c'z + d') = 0, \; \lambda, \mu \in \mathbb{R}
    \]



\subsection{Rette complanari e rette sghembe}
    Date due rette $r$ ed $s$, queste si dicono complanari se e solo se sono parallele o incidenti. Si dicono invece sghembe se e solo se non sono parallele e non sono complanari. 

    \vspace{10pt}
    Dati due punti $P$ giacente su $r$ e $Q$ giacente su $s$, la distanza tra due rette sghembe è data da:
    \[
        d(r, s) = \frac{|\mathbf{v} \times \mathbf{w} \cdot (P - Q)|}{||\mathbf{v} \times \mathbf{w}||}
    \]
    Dove $\mathbf{v}, \mathbf{w}$ sono i vettori direzione di $r$ ed $s$.



\subsection{Distanza punto-piano}
    Dato un punto $P$ ed un piano $\pi : ax + by + cz + d = 0$ la distanza tra di loro è data da: 
    \[
        d(P, \pi) = \frac{|(a, b, c) \cdot (P - P_0)|}{||(a, b, c)||} = \frac{|a(x - x_0) + b(y - y_0) + c(z - z_0)|}{\sqrt{a^2 + b^2 + c^2}} = \frac{|ax + by + cz + d|}{\sqrt(a^2 + b^2 + c^2)}
    \]
    



\end{document}